\chapter{S2：全温与温度补偿验证}
\section{阶段定位与设备边界}
本阶段在 S1 常温基线的基础上评价温度对偏置、标度、角度和算法输出的影响。\textbf{不得将现有 2TD-600 当成具温箱的设备：其说明书规定工作环境 $25\pm5\,{}^\circ\mathrm C$，若温箱与参考转台无法安全兼容，应重新设计设备组合。}
\section{S2 测试项目总表}
\small
\begin{longtable}{C{18mm}p{46mm}p{61mm}p{22mm}}
\toprule 编号 & 测试项 & 测试目标／方法提示 & 级别／条件\\\midrule\endfirsthead
\toprule 编号 & 测试项 & 测试目标／方法提示 & 级别／条件\\\midrule\endhead
S2-00 & 温控设备、参考与安全核查 & 温箱/线缆/转台是否可组合使用；温度探头与时间同步 & 前置/待设备核实\\
S2-01 & 温度稳态和热平衡验证 & 多测点热平衡、芯片与空气温度差 & P0\\
S2-02 & ACC 零偏温度残差 & 三轴零偏随温度变化、冻参基线 & P0\\
S2-03 & GYRO 零偏温度残差 & 三轴静止零速曲线和补偿残余 & P0\\
S2-04 & ACC 比例因子温度变化 & 多个参考姿态下矩阵残差随温度变化 & P1\\
S2-05 & GYRO 比例因子温度变化 & 不同温度多速率；需可溯源速率 & P1/有条件\\
S2-06 & 交叉轴/非正交参数温漂 & 多位置/速率联合检查矩阵项稳定性 & P1\\
S2-07 & 内置温度传感器准确性 & 芯片读数对外部参考温度、响应延迟 & P1\\
S2-08 & 冷/热机预热和开机重复性 & 温度极值启动、归零时间及重复性 & P0\\
S2-09 & 升降温迟滞 & 相同温度点升温/降温路径比较 & P0\\
S2-10 & 温变速率和热梯度瞬态 & 动态热过程与静态稳态偏差区分 & P1\\
S2-11 & 温度循环后的参数复现 & 多个循环前后静态和速率残差 & P1\\
S2-12 & 温度补偿前后独立对比 & 独立于补偿拟合的数据、误差改善与异常 & P0\\
S2-13 & 全温静态倾角精度 & 不同温度下单轴/复合姿态误差 & P0\\
S2-14 & 全温噪声与 Allan 复核 & 选定温度点复用 CT-12A 分析法 & P1\\
S2-15 & 全温动态角度及阶段报告 & 具动态角基准时才验证动态输出；全温归档 & P1/有条件\\
\bottomrule
\end{longtable}
\normalsize
\section{逐项方案编写接口（后续版本）}
本阶段\textbf{测试编号已预留，详细实施方案尚未完成}。每个编号在正式执行前，应按 S1 章节的模板补齐：目的、仪器和参考不确定度、进入条件、工况矩阵、台架操作、原始数据字段、计算公式、合格判据、异常处理和结果归档。不得因本表列有测试名称就自动判定已经执行或通过。
\begin{warningbox}温度阶段的具体温度点、驻留时间、热梯度、升降温速率与绝对精度限值均须依据产品工作温度和温控设备条件确定。不要把 2FDT600WK 的温控功能推定为 2TD-600 能力。\end{warningbox}
